% Part: computability
% Chapter: computability-theory
% Section: total

\documentclass[../../../include/open-logic-section]{subfiles}

\begin{document}

\olfileid{cmp}{thy}{tot}
\olsection{సర్వనిర్వచితత్వం అనిర్ణయనీయం}

ఏదో ఒకటి గణనీయం కాదని చూపడానికి $s$-$m$-$n$ సిద్ధాంతాన్ని వాడే
మరొక ఉదాహరణను పరిశీలిద్దాం. $\fn{Tot}$ అనేది సర్వనిర్వచిత గణనీయ
ప్రమేయాల సూచికల సమితి అనుకుందాం; అంటే
\[
\fn{Tot}=\Setabs{x}{\text{ప్రతి $y$కు $\cfind{x}(y)\fdefined$}}.
\]

\begin{prop}
\ollabel{prop:total}
$\fn{Tot}$ గణనీయం కాదు.
\end{prop}

\begin{proof}
$\fn{Tot}$ గణనీయం కాదని చూడటానికి, $K$ దానికి తగ్గించదగినదని చూపితే
చాలు. $h(x,y)$ను ఇలా నిర్వచించండి:
\[
h(x,y) \simeq
\begin{cases}
0 & \text{$x\in K$ అయితే} \\
\fundefined & \text{ఇతర సందర్భాల్లో.}
\end{cases}
\]
$h(x,y)$ అసలు $y$పై ఆధారపడదని గమనించండి. $h$ పాక్షిక గణనీయమని
చూడటం కష్టం కాదు: $h$ను గణించడానికి, ఇన్‌పుట్‌లు $x,y$పై ముందుగా ఇన్‌పుట్~$x$తో
ప్రమేయం~$\cfind{x}$ను అనుకరిస్తాం; ఆ గణన ఆగితే
$h(x,y)$ విలువ~$0$ను ఇచ్చి ఆగుతుంది. కాబట్టి $h(x,y)$ అనేది
$\Zero(\umin{s}{T(x,x,s)})$ మాత్రమే; ఇక్కడ $\Zero$ స్థిర శూన్య
ప్రమేయం.
\sourcecorrection{OLTECOMTHY-016}{ఈ గణనా వివరణలో ఆంగ్ల మూలం ``then we compute''కు బదులు ``the we compute'' అనే అక్షరదోషంతో రాసింది. ఉద్దేశించిన ప్రక్రియను సరిగా అనువదించాం.}

$s$-$m$-$n$ సిద్ధాంతం ప్రకారం, ప్రతి $x$,~$y$లకు
\[
\cfind{k(x)}(y) =
\begin{cases}
0 & \text{$x\in K$ అయితే} \\
\fundefined & \text{ఇతర సందర్భాల్లో}
\end{cases}
\]
అయ్యే ఆదిమ పునరావృత్త ప్రమేయం~$k(x)$ ఒకటి ఉంది. అందువల్ల $x\in K$
అయితే $\cfind{k(x)}$ సర్వనిర్వచితం; లేకపోతే అది ఎక్కడా నిర్వచితం
కాదు. కనుక $k$, $K$ను~$\fn{Tot}$కు తగ్గిస్తుంది.
\end{proof}

\begin{digress}
నిజానికి $\fn{Tot}$ గణనీయంగా లెక్కించదగినది కూడా కాదు---దాని
సంక్లిష్టత ``అంకగణిత శ్రేణీకరణ''లో మరింత పైస్థాయిలో ఉంటుంది. అయితే
ఈ బలమైన ఫలితం గురించి ఇక్కడ ఆందోళన పడనవసరం లేదు.
\end{digress}

\end{document}
